% Dot batch 78; source SHA-256: df60dc4aa19b733f13bd11225b51651f4d51d22435b07c8be116aecb3acdcdbe
% Archived source: dot_updates_20261124/batches/batch078/EGA第五卷_线性系与关联除子.tex
% Original body lines: 26--88; integration map: integration/late_integration.json
\begingroup
\def\newscope{V-L}
\def\Div{\operatorname{Div}}
\def\Supp{\operatorname{Supp}}
\def\Ass{\operatorname{Ass}}
\def\Spec{\operatorname{Spec}}
\def\PP{\mathbf P}
\def\OO{\mathcal O}
沿用上文：\(X\to S\) 平坦且局部有限表示，\(D\) 是由 \(T\) 参数化的 \(X/S\) 上的除子族，\(Z\) 是固定点集。
现在设 \(T\) 是射影丛 \(Q=\PP(F)\)，\(F\) 是有限型局部自由模；以下在 \(Q\) 的纤维非空处讨论。
记
\[
 P=\PP(F^\vee),\qquad Q=\PP(F)=P^\vee,\qquad
 P(n)=\PP\bigl((\operatorname{Sym}^n F)^\vee\bigr)\quad(n\geq0).
\]
若 \(F\) 的秩处处至少为二，则 \(\Div_{Q/S}\) 由各个 \(P(n)\) 的不交并表出，转置给出一个态射
\[
 X\setminus Z\longrightarrow\coprod_{n\geq0}P(n).
\]
如果该态射经 \(P(n)\) 因子分解，就称 \(D\) 是\emph{\(n\) 次的}。
此时若 \(X\ne\varnothing\)，则 \(X\setminus Z\ne\varnothing\)，因而这个 \(n\) 由 \(D\) 唯一确定。
不过，要定义给定的次数，严格说来不必用到上面的可表性结果，只须定义典范单态射
\[
 P(n)\longrightarrow\Div_{P^\vee/S}.
\]
于是也可直接要求，转置态射经这一个单态射因子分解。\cmt{L18-078-01}{法文物理129页和英译物理98页都是先说由“各 \(P(n)\) 的和”表出，却把一般态射的目标径写为 \(P(n)\)；应先写不交并。法、英稿还把对偶放在 \(F\) 上。按前文2.2--2.4与商约定，次数 \(n\) 的方程给出 \(\operatorname{Sym}^n F\) 中的可逆直和因子，取对偶才得到 \((\operatorname{Sym}^n F)^\vee\) 的可逆商。底上 \(n!\) 不可逆时，不能一般地用 \(\operatorname{Sym}^n(F^\vee)\) 替代这个对偶；\(n=1\) 不受影响。其各分支为开闭，故次数在定义域上局部常值，不必全局相同。原文还漏列了此可表性分解的秩条件：秩一时 \(Q=S\)，其相对有效除子只能为空；各个 \(P(n)\simeq S\) 给出的是同一个除子函子点，并非彼此不交。此时仍可用 \(n=1\) 的典范单态射定义一次线性系。秩零处不在参数纤维非空的讨论内。}

所谓\emph{由射影丛 \(Q=P^\vee\) 参数化的 \(X/S\) 上的除子线性系}，
是由 \(Q\) 参数化且为一次的除子族，即其转置确定一个 \(S\)-态射
\[
 f:X\setminus Z\longrightarrow P.
\]
因而，当 \(P^\vee\) 的纤维非空时，这样一个线性系典范地给出由 \(X\) 到射影丛的有理映射；更准确地说，是一个\emph{相对于 \(S\) 的伪态射}。\cmt{L18-078-02}{法稿物理129页的手写补充为参数纤维非空及相对伪态射。英译脚注76把前一处标为不清；法稿仍可据以辨认。此处保留作者术语，不把相对概形稠密直接换成任意基底上的通常有理映射定义。}
按构造，
\[
 D\big|_{(X\setminus Z)\times_S P^\vee}
 =
 (f\times_S\mathrm{id}_{P^\vee})^{-1}(H),
\]
其中 \(H\) 是 \(P\times_S P^\vee\) 上的典范除子，即\emph{关联除子}。
因此知道 \(f\) 就至少可以重建 \(D\) 在 \(X\setminus Z\) 上诱导的除子族。
特别地，若原族无固定点，其关联态射 \(f:X\to P\) 就决定整个除子族。

无固定点的这类线性系，与使
\((f\times_S\mathrm{id}_{P^\vee})^{-1}(H)\)
成为 \(X\times_S P^\vee\) 上相对于 \(P^\vee\) 的除子的态射 \(f:X\to P\)，
构成一一对应。这个条件可以逐纤维检验，得到：\cmt{L18-078-03}{法文物理130页明确写“相对于 \(P^\vee\)”。
英译物理99页的排印次序易混淆，此处依照法稿恢复。一次转置从相对于 \(X\) 的一侧给定；要得到 \(X/S\) 上由 \(P^\vee\) 参数化的除子族，还须检查相对于 \(P^\vee\) 的另一侧。}

\egas{3.1}{命题}

由 \(P^\vee\) 参数化的 \(X/S\) 上无固定点的除子线性系，与满足下列条件的 \(S\)-态射 \(f:X\to P\) 一一对应：
对每个 \(s\in S\)，取剩余域 \(k(s)\) 的一个代数闭包 \(\bar k\)，并令 \(f_{\bar k}:X_{\bar k}\to P_{\bar k}\) 为相应的基变换。
对 \(X_{\bar k}\) 的每个关联素循环 \(X'\)，
\(f_{\bar k}(X')\) 都不包含于 \(P_{\bar k}\) 的任何超平面。
如果 \(X\) 的纤维几何整，这一条件只须说：
\(f_{\bar k}(X_{\bar k})\) 不包含于 \(P_{\bar k}\) 的任何超平面。\cmt{L18-078-04}{命题编号依手改后的法稿和英译。这里的关联素循环是关联点闭包所给的整闭子概形，不只包括极小关联点对应的不可约分支。遇到嵌入关联点，也必须检查。}

一般地，即当 \(Z\ne\varnothing\) 时，\(f\) 不再一定能决定除子族。
最简单的情形是 \(P^\vee=S\)，即相对维数为零。给定一个由 \(S\) 参数化的
\(X/S\) 上的除子线性系，就等于给定 \(X\) 上相对于 \(S\) 的一个有效 Cartier 除子 \(D\)。
关联态射是投影 \(X\setminus\Supp D\to S\)。
即使同时知道这个态射的定义开集，也不能据此确定 \(D\)。\cmt{L18-078-05}{例如，取 \(X=\mathbb A^1_k\)、\(S=\Spec k\)，原点及其二重除子有相同的支集，给出同一个补开集及同一个投影。这里 \(P^\vee=S\) 是秩一的情形；\(P(1)\to\Div_{P^\vee/S}\) 对应空除子，故转置的次数一条件对这个补开集自动成立。}
在此情形，若为简便而设 \(X\) 约化，把 \(f\) 看成从 \(X\) 到 \(P=S\) 的有理映射，
其最大定义开集不是 \(X\setminus\Supp D\)，而是 \(X\) 本身。

为排除这种现象，以下限制为“\emph{无固定分支}”的除子线性系。
一般地，若 \(S=\Spec k\)，对一个由 \(T\) 参数化的 \(X/S\) 上的除子族（未必线性），
固定点集 \(Z\) 中任一在 \(X\) 内余维为一的不可约分支，称为一个\emph{固定分支}。
若无这样的分支，即 \(\operatorname{codim}_X Z\geq2\)，就称这个族无固定分支。
对任意底概形 \(S\)，逐纤维使用这一定义即可；无固定分支的性质在基变换下保持。\cmt{L18-078-06}{此处是上文使固定点集为闭集的语境。纤维上，参数空间非空时，\(Z\) 包含于任选一个成员的支集，因而不含任何余维零分支。可在域扩张后作这一检查。域扩张保持有限型几何纤维中闭子集分支的余维；固定点集又与基变换相容。无固定分支不等于没有固定点：允许有余维至少二的固定点。}
\par\endgroup
